\documentclass{article}
\usepackage[utf8]{inputenc}
\usepackage[spanish]{babel}
\usepackage{listings}
\usepackage{xcolor}
\usepackage{geometry}
\geometry{a4paper, margin=1in}

\% Configuración para código C++
\definecolor{codegreen}{rgb}{0,0.6,0}        
\definecolor{codegray}{rgb}{0.5,0.5,0.5}     
\definecolor{codepurple}{rgb}{0.58,0,0.82}   
\definecolor{backcolour}{rgb}{0.95,0.95,0.92}

\lstset{
    backgroundcolor=\color{backcolour},      
    commentstyle=\color{codegreen},
    keywordstyle=\color{magenta},
    numberstyle=\tiny\color{codegray},       
    stringstyle=\color{codepurple},
    basicstyle=\ttfamily\small,
    breakatwhitespace=false,
    breaklines=true,
    captionpos=b,
    keepspaces=true,
    numbers=left,
    numbersep=5pt,
    showspaces=false,
    showstringspaces=false,
    showtabs=false,
    tabsize=2,
    language=C++,
    inputencoding=utf8,
    extendedchars=true,
    literate={á}{{\'a}}1 {é}{{\'e}}1 {í}{{\'i}}1 {ó}{{\'o}}1 {ú}{{\'u}}1 {ñ}{{\~n}}1 {¿}{{\textquestiondown}}1 {¡}{{\textexclamdown}}1
}

\% Macros para las secciones de los apuntes
\newcommand{\quees}[1]{
    \section*{🤔 ¿QUÉ ES?}
    \hrulefill \\
    \#1
}

\newcommand{\dichodeotraforma}[1]{
    \section*{💬 DICHO DE OTRA FORMA}
    \hrulefill \\
    \#1
}

\newcommand{\diagrama}[1]{
    \section*{🎨 DIAGRAMA}
    \hrulefill \\
    \#1
}

\newcommand{\comofunciona}[1]{
    \section*{⚙️ ¿CÓMO FUNCIONA?}
    \hrulefill \\
    \#1
}

\newcommand{\complejidad}[1]{
    \section*{📱 COMPLEJIDAD (Big-O)}
    \hrulefill \\
    \#1
}

\newcommand{\entrevistas}[1]{
    \section*{🎯 EN ENTREVISTAS}
    \hrulefill \\
    \#1
}

\newcommand{\resumen}[1]{
    \section*{📋 RESUMEN RÁPIDO}
    \hrulefill \\
    \#1
}

\begin{document}

\title{Heaps (Montículos)}
\author{Aldo Zepeda Montoya}
\date{\today}
\maketitle

\subsection*{CATEGORÍA: 02\_Estructuras\_No\_Lineales}

\quees{
Un torneo de fútbol ⚽ donde el mejor equipo SIEMPRE llega a la 
final (está en la cima). 

Si el mejor equipo pierde o se retira (lo sacas), el siguiente mejor
sube automáticamente a ocupar su lugar. No sabes exactamente quién es
el tercer o cuarto mejor sin hacer más comparaciones, pero SIEMPRE 
tienes al \#1 disponible de inmediato.

Es un árbol binario especial donde la raíz siempre tiene el valor 
máximo (Max-Heap) o mínimo (Min-Heap) de toda la estructura.
}

\dichodeotraforma{
Un Heap es un árbol binario completo (lleno de arriba a abajo, de
izquierda a derecha) que cumple la "Propiedad de Heap":
- Max-Heap: El padre siempre es mayor o igual que sus hijos.
- Min-Heap: El padre siempre es menor o igual que sus hijos.

A diferencia del BST, no hay orden entre hermanos. Lo único que 
importa es la relación Padre-Hijo. Por su estructura compacta, 
se suele implementar eficientemente usando un Array.
}

\diagrama{\begin{verbatim}
Max-Heap Visual:
          [90]
         /    \
       [80]  [75]
       / \   / \
     [60][70][50][55]

Implementado como Array (¡Truco de índices!):
Idx:  [0] [1] [2] [3] [4] [5] [6]
Val: [90][80][75][60][70][50][55]

Fórmulas mágicas para el índice `i`:
- Parent: (i - 1) / 2
- Left Child: 2 * i + 1
- Right Child: 2 * i + 2
\end{verbatim}}

\begin{lstlisting}[language=C++]
1. Insert: Agregas el nuevo valor al final del array (la hoja más a
   la izquierda disponible). Luego lo comparas con su padre y lo
   subes si es necesario (Bubble Up / Heapify Up). O(log n).
2. Extract Max/Min: Sacas la raíz (arr[0]). Mueves el último elemento
   a la raíz y lo vas bajando comparándolo con sus hijos hasta que
   esté en su lugar (Bubble Down / Heapify Down). O(log n).
3. Peek: Miras arr[0]. O(1).
4. Heapify (Construir Heap): Convertir un array desordenado en un
   Heap. O(n).

📝 PSEUDOCÓDIGO
──────────────────────────────────────────────────────────────

// Heapify Up (al insertar)
función subir(i):
    MIENTRAS i > 0 Y arr[padre(i)] < arr[i]:
        intercambiar(arr[padre(i)], arr[i])
        i = padre(i)

// Heapify Down (al extraer)
función bajar(i):
    mayor = i
    SI hijoIzquierdo(i) existe Y arr[hijoIzquierdo(i)] > arr[mayor]:
        mayor = hijoIzquierdo(i)
    SI hijoDerecho(i) existe Y arr[hijoDerecho(i)] > arr[mayor]:
        mayor = hijoDerecho(i)
    SI mayor != i:
        intercambiar(arr[i], arr[mayor])
        bajar(mayor)

💻 CÓDIGO C++
──────────────────────────────────────────────────────────────

#include <iostream>
#include <vector>
#include <queue>    // STL priority_queue es un Heap
using namespace std;

int main() {
    // 1. Max-Heap (por defecto)
    priority_queue<int> maxHeap;
    
    maxHeap.push(10);
    maxHeap.push(30);
    maxHeap.push(20);

    cout << "Máximo en la cima: " << maxHeap.top() << endl; // 30

    maxHeap.pop(); // Quita el 30
    cout << "Nuevo máximo: " << maxHeap.top() << endl;    // 20

    // 2. Min-Heap
    priority_queue<int, vector<int>, greater<int>> minHeap;
    minHeap.push(10);
    minHeap.push(30);
    minHeap.push(20);

    cout << "Mínimo en la cima: " << minHeap.top() << endl; // 10

    return 0;
}

⏱️ COMPLEJIDAD (Big-O)
──────────────────────────────────────────────────────────────

  Operación        │ Tiempo    │ Razón
  ─────────────────┼───────────┼──────────────────────────────
  Insert           │ O(log n)  │ Altura del árbol
  Extract Max/Min  │ O(log n)  │ Altura del árbol
  Peek Max/Min     │ O(1)      │ Es la raíz (arr[0])
  Heapify (Build)  │ O(n)      │ Análisis matemático avanzado
  Search           │ O(n)      │ No es un BST, hay que buscar todo
\end{lstlisting}

\begin{lstlisting}[language=C++]
✅ "Top K Elements": Si te piden los K elementos más grandes, usa un
   Min-Heap de tamaño K. Si te piden los más chicos, usa un Max-Heap.

💡 Priority Queue = Heap: En el 99% de los casos, cuando necesites un 
   Heap en una entrevista, usa `std::priority_queue`. No lo programes
   desde cero a menos que te lo pidan explícitamente.

⚠️ No es un BST: Recuerda que un Heap NO está ordenado de izquierda a
   derecha. Solo garantiza la relación Padre > Hijo.

🔑 Heapsort: Puedes ordenar un array metiendo todo a un Heap y 
   sacándolo uno por uno. O(n log n).
\end{lstlisting}

\resumen{
• Árbol binario completo donde el jefe (raíz) es el mayor o menor.
• Implementado como Array para máxima velocidad y ahorro de memoria.
• Insertar y sacar el máximo/mínimo es O(log n).
• Ver el máximo/mínimo es instantáneo O(1).
• Usos: Algoritmo de Dijkstra, colas de prioridad, problemas de "Top K".
• Heapify (construir desde array) es sorprendentemente O(n).
════════════════════════════════════════════════════════════════
}

\end{document}